\section{The canonical Rees degeneration and normal cone}
\label{sec:rees-specialization}

We pass from the associated graded algebra to the specialization
mechanism that produces it. This places every projection rank in one
special fibre over the whole Schubert divisor, while the construction
itself is attached to the abstract nonreduced scheme.

\begin{proof}[Proof of Theorem~\ref{thm:canonical-rees-specialization}]
The nilradical \(\widehat{\mathcal N}_R\) is intrinsic, coherent, and
nilpotent. On an affine open
\(\operatorname{Spec}A\subset\widehat D_R\), write \(N\subset A\)
for its nilradical and set
\[
 \mathscr R_N=A[t,Nt^{-1}]\subset A[t,t^{-1}].
\]
This is the subalgebra generated by \(A[t]\) and the elements
\(nt^{-1}\), \(n\in N\). Formation of the nilradical and its powers
commutes with localization, so these affine algebras glue.

Multiplication by \(t\) is injective on \(\mathscr R_N\), since it is
a subring of \(A[t,t^{-1}]\). Thus every stalk is torsion-free as a
\(\C[t]\)-module. Since \(\C[t]\) is a principal ideal domain,
torsion-free \(\C[t]\)-modules are flat. Hence the resulting scheme is
flat over \(\A^1\).

After inverting \(t\),
\[
 \mathscr R_N[t^{-1}]=A[t,t^{-1}],
\]
so the restriction over \(\mathbb G_m\) is
\(\widehat D_R\times\mathbb G_m\). For the special fibre, the
extended-Rees calculation gives
\[
 \mathscr R_N/(t)
 \simeq
 \bigoplus_{j\ge0}N^j/N^{j+1}
 =
 \operatorname{gr}_N A.
\]
Indeed the degree-\(j\) class is represented by \(nt^{-j}\),
\(n\in N^j\), and vanishes modulo \(t\) precisely when
\(n\in N^{j+1}\). These identifications commute with localization.
Consequently
\[
 \mathfrak D_{R,0}
 \simeq
 \operatorname{Spec}_{\widehat\Delta_R}
 \operatorname{gr}_{\widehat{\mathcal N}_R}
 \mathcal O_{\widehat D_R}.
\]
Because \(\widehat{\mathcal N}_R\) is the ideal of
\(\widehat\Delta_R=(\widehat D_R)_{\mathrm{red}}\) in
\(\widehat D_R\), the right-hand side is the normal cone
\(C_{\widehat\Delta_R/\widehat D_R}\).

Theorem~\ref{thm:global-residual-colon} identifies each graded piece
with
\[
 \mathcal L^{\otimes j}\otimes\mathcal O_{W_j},
\]
and its multiplication with multiplication of successive powers of
the Schubert ideal, equivalently with the compatible maps in the
colon tower. An abstract scheme isomorphism preserves the nilradical,
so it preserves the extended Rees algebra, the family, and its special
fibre. This proves functoriality.
\end{proof}

\begin{corollary}[One specialization across the first boundary]
\label{cor:three-boundary-regimes}
The canonical special fibre restricts as follows.
\begin{enumerate}
\item On the adjacent projection-corank-one locus it truncates after
degree one, whose support is \(E_R\).
\item On the mixed-regular rank-two locus it is
\[
 \mathcal O_{\widehat\Delta_R}
 \oplus
 \bigl(\mathcal L\otimes\mathcal O_{\mathcal C_H}\bigr)
 \oplus
 \bigl(\mathcal L^{\otimes2}\otimes\mathcal O_{\Sigma_R}\bigr).
\]
Contact collisions alter the Cartier multiplicities in
\(\mathcal C_H\) but leave the vertex-primary factor present.
\item At a transverse decomposable mixed-kernel wall, the degree-two
support is
\[
 V(c,d,a^2,ab,b^2),
\]
a length-three thickening of the rank-two vertex.
\end{enumerate}
These are restrictions of one canonical normal cone.
\end{corollary}

\begin{proof}
The first assertion is the corank-one normal form \(t(t,f)\). The
second follows from Theorems~\ref{thm:saturated-contact} and
\ref{thm:graded-nilpotent}; the third is
Theorem~\ref{thm:decomposable-kernel-wall}. The preceding theorem
places all three in the same special fibre.
\end{proof}

\begin{remark}[Presentation versus structure]
The residual block identity of
Theorem~\ref{thm:all-corank-presentation} is a presentation on a
chosen Schur chart. By contrast, \(\mathfrak D_R\to\A^1\) is built
from the intrinsic nilradical. Its special fibre records every
annihilator layer and every multiplication map between successive
layers. Thus the Rees family supplies the global specialization
theorem requested by the higher-corank objection; the block formula
remains a tool for computing particular strata of that object.
\end{remark}
